\chapter*{Appendix}
\addcontentsline{toc}{chapter}{Appendix}


\section*{Additional Material}
\addcontentsline{toc}{section}{Additional Material}


\begin{table}[H]
	\renewcommand{\arraystretch}{1.5}
	\begin{center}
		\begin{tabularx}{\textwidth}{|p{2.1cm}|Y|Y|}
			\hline
			& \textbf{Proof of $\mathbb{P}(a(T)\leq 0)=\alpha$}
			& \textbf{Proof of $\mathbb{P}(b(T)\geq 0)=\alpha$}
			\\
			\hline
			\textbf{Preconditions}
			&
			\begin{itemize}[leftmargin=*,nosep]
				\item $a$ is left-continuous.
				\item $\mathbb{P}(\theta(u,t)\geq a(t))=\alpha
				\quad \text{for all } t$
			\end{itemize}
			&
			\begin{itemize}[leftmargin=*,nosep]
				\item $b$ is right-continuous.
				\item $\mathbb{P}(\theta(u,t)\leq b(t))=\alpha
				\quad \text{for all } t$
			\end{itemize}
			\\
			\hline
			\textbf{Step (i)}
			&
			Define
			\(
			F^\theta(t):=
			\mathbb{P}(\theta(u,v)\leq t)
			\)
			\begin{itemize}[nosep]
				\item[$\rightarrow$] right-continuous in $t$
				\item[$\rightarrow$] monotonically decreasing in $\theta$
			\end{itemize}
			&
			Define
			\(
			G^\theta(t):=
			\mathbb{P}(\theta(u,v)\geq t)
			\)
			\begin{itemize}[nosep]
				\item[$\rightarrow$] left-continuous in $t$
				\item[$\rightarrow$] monotonically increasing in $\theta$
			\end{itemize}
			\\
			\hline
			\textbf{Step (ii)}
			&
			Show:
			\(
			F^{a(t)}(t)=\alpha
			\)
			for all $t$
			&
			Show:
			\(
			G^{b(t)}(t)=\alpha
			\)
			for all $t$
			\\
			\hline
			\textbf{Step (iii)}
			&
			Define
			$S:=\{t:a(t)\leq \theta\}$
			
			Show that
			\[
			\begin{aligned}
				F^\theta(t)>\alpha &\implies t\notin S,\\
				F^\theta(t)<\alpha &\implies t\in S.
			\end{aligned}
			\]
			Define
			\[
			\begin{aligned}
				q_1&:=\inf\{t:F^\theta(t)\geq\alpha\},\\
				q_2&:=\inf\{t:F^\theta(t)>\alpha\}.
			\end{aligned}
			\]
			Then $q_1\leq q_2$.
			Split $S=S_1\cup S_2\cup S_3$, with
			\[
			\begin{aligned}
				S_1&=S\cap(-\infty,q_1],\\
				S_2&=S\cap(q_1,q_2),\\
				S_3&=S\cap[q_2,\infty).
			\end{aligned}
			\]
			&
			Define
			$S':=\{t:b(t)\geq \theta\}$
			
			Show that
			\[
			\begin{aligned}
				G^\theta(t)>\alpha &\implies t\notin S',\\
				G^\theta(t)<\alpha &\implies t\in S'.
			\end{aligned}
			\]
			Define
			\[
			\begin{aligned}
				q_1'&:=\sup\{t:G^\theta(t)\geq\alpha\},\\
				q_2'&:=\sup\{t:G^\theta(t)>\alpha\}.
			\end{aligned}
			\]
			Then $q_1'\geq q_2'$.
			Split $S'=S_1'\cup S_2'\cup S_3'$, with
			\[
			\begin{aligned}
				S_1'&=S'\cap[q_1',\infty),\\
				S_2'&=S'\cap(q_2',q_1'),\\
				S_3'&=S'\cap(-\infty,q_2'].
			\end{aligned}
			\]
			\\
			\hline
			\textbf{Step (iv)}
			\\[-2mm]
			&
			Observe that $\{T\leq q_1\}=\{T\in S_1\}$
			
			It follows that
			\[
			\begin{aligned}
				\mathbb{P}(T\in S_1)
				&=\mathbb{P}(T\leq q_1)\\
				&=F^\theta(q_1)\\
				&=\alpha.
			\end{aligned}
			\]
			&
			Observe that $\{T\geq q_1'\}=\{T\in S_1'\}$
			
			It follows that
			\[
			\begin{aligned}
				\mathbb{P}(T\in S_1')
				&=\mathbb{P}(T\geq q_1')\\
				&=G^\theta(q_1')\\
				&=\alpha.
			\end{aligned}
			\]
			\\
			\hline
			\textbf{Step (v)}
			\\[-15mm]
			&
			\[
			\begin{aligned}
				\mathbb{P}(T\in S_2)
				&\leq
				\lim_{t\uparrow q_2}F^\theta(t)
				-F^\theta(q_1)\\
				&=\alpha-\alpha\\
				&=0.
			\end{aligned}
			\]
			&
			\[
			\begin{aligned}
				\mathbb{P}(T\in S_2')
				&\leq
				\lim_{t\downarrow q_2'}G^\theta(t)
				-G^\theta(q_1')\\
				&=\alpha-\alpha\\
				&=0.
			\end{aligned}
			\]
			\\
			\hline
			\textbf{Step (vi)}
			\\[-6mm]
			&
			Show: $t>q_2\implies t\notin S$
			
			\begin{itemize}[nosep]
				\item[$\rightarrow$] $S_3\subseteq\{q_2\}$
				\item[$\rightarrow$] $\mathbb{P}(T\in S_3)=0$
			\end{itemize}
			&
			Show: $t<q_2'\implies t\notin S'$
			
			\begin{itemize}[nosep]
				\item[$\rightarrow$] $S_3'\subseteq\{q_2'\}$
				\item[$\rightarrow$] $\mathbb{P}(T\in S_3')=0$
			\end{itemize}
			\\
			\hline
			\textbf{Step (vii)}
			\\[-2mm]
			& Calculate $\mathbb{P}(T \in S)$
			& Calculate $\mathbb{P}(T \in S')$
			\\
			\hline
		\end{tabularx}
		
		\caption{Summary for the proof of exactness in Corollary~\ref{corollary:Corollary 2 Lindqvist}}
		\label{tab:comparison a b}
	\end{center}
\end{table}



\newpage
\section*{Calculations}
\addcontentsline{toc}{section}{Calculations}

\subsection{Distributions of $\bar \theta_{\text{full}}(\mathbf u, 2 / 3)$ and $\bar \theta_{\text{fid}}(\mathbf u, 2 / 3)$}



In the following, we will obtain the distribution of $\hat \theta_{\text{fid}} (\mathbf U, 2/3) = 2 U_{1:2} \; | \; 2 U_{1:2} \leq U_{2:2}$. 


For easier notation, let us define $M = U_{1:2}$.
We can write
\[
U_{2:2} = M + E
\]
with $E \sim \text{Exp}(1)$ because of the memoryless property of the exponential distribution.
The minimum $M = \min(U_1, U_2)$ is the time of the first event. 
At this time, one of the two variables is "used up". 
The other variable has not happened yet. 
By the memoryless property, the remaining waiting time until it happens has the same $\text{Exp}(1)$ distribution as at the start, and it is independent of $M$.
Therefore, the difference $E = U_{2:2} - U_{1:2}$ is $\text{Exp}(1)$-distributed and independent of $M$.

We can also show this rigerously by applying the transformation theorem on the vector $(U_{1:2}, U_{2:2})$ and the function $T(x,y) = (x, y-x)$.
Integrating out the first component gives $U_{2:2} - U_{1:2} = E \sim \text{Exp}(1)$.

Next, note that
\[
2 U_{1:2} \leq U_{2:2}
\iff 2M \leq M + E
\iff M \leq E.
\]
So we can use this notation to describe the fiducial distribution as
\[
\hat{\theta}_{\text{fid}}(\mathbf{U}, 2/3) \sim 2M \;\big|\; M \leq E,
\]
where $M \sim \text{Exp}(2)$ and $E \sim \text{Exp}(1)$ are independent.


\begin{align*}
	\mathbb{P} \left( \hat{\theta}_{\text{fid}}(\mathbf{U}, 2/3) \leq x  \right)
	=
	\mathbb{P} \left( 2M \leq x \;\big|\; M \leq E  \right) \\
	=
	\frac{
		\mathbb{P} \left( 2M \leq x , M \leq E  \right)
	}{
		\mathbb{P} \left( M \leq E  \right)
	}
\end{align*}
The denominator is
\begin{align*}
	\mathbb{P} \left( M \leq E  \right) 
	&=
	\int_0^\infty \mathbb{P} (m \leq E \;\big|\; M = m) \; dM \\
	&=
	\int_0^\infty \mathbb{P} (m \leq E \;\big|\; M = m) f_M(m) \; dm \\
	&=
	\int_0^\infty \exp(-m) \cdot 2 \exp(-2m) \; dm \\
	&=
	\frac{2}{3}
\end{align*}
and for $x \geq 0$, the numerator is
\begin{align*}
	\mathbb{P} \left( 2M \leq x , M \leq E  \right)
	&=
	\mathbb{P} \left( M \leq \tfrac{x}{2} , M \leq E  \right) \\
	&=
	\int_0^{x/2} \mathbb{P} (m \leq E) \, f_M(m) \; dm \\
	&=
	\int_0^{x/2} \exp(-m) \cdot 2 \exp(-2m) \; dm \\
	&=
	\frac{2}{3} \left( 1 - \exp\!\left( -\frac{3x}{2} \right) \right).
\end{align*}
Putting both parts together, we get for $x \geq 0$
\begin{align*}
	\mathbb{P} \left( \hat{\theta}_{\text{fid}}(\mathbf{U}, 2/3) \leq x  \right)
	=
	\frac{\frac{2}{3} \left( 1 - \exp\left( -\frac{3x}{2} \right) \right)}{\frac{2}{3}}
	=
	1 - \exp\!\left( -\frac{3x}{2} \right).
\end{align*}
For $x < 0$ the probability is $0$.
This is the CDF of an exponential distribution with rate $3/2$, so
\[
\hat{\theta}_{\text{fid}}(\mathbf{U}, 2/3) \sim \text{Exp}\!\left( \tfrac{3}{2} \right).
\]


Next, we will calculate the CDF and PDF of \(\bar{\theta}_{\text{full}}(\mathbf{U},2/3)\).
It holds that
\[
\bar{\theta}_{\text{full}}(\mathbf{U},2/3)
=
\min\left(2U_{1:2},U_{2:2}\right)
=
\min\left(2M,M+E\right).
\]
The event \(\{\bar{\theta}_{\text{full}}(\mathbf{U},2/3)>x\}\) is equivalent to
\(\{M>x/2\}\cap\{E>x-M\}\). Thus, using Fubini
\begin{align*}
	\mathbb{P}\left(\bar{\theta}_{\text{full}}(\mathbf{U},2/3)>x\right)
	&=
	\mathbb{P}(M>x/2,E>x-M)\\
	&=
	\int_{x/2}^{\infty}
	\mathbb{P}(E>x-m)f_M(m)\,\mathrm{d}m.
\end{align*}
Moreover,
\[
\mathbb{P}(E>y)
=
\begin{cases}
	\exp(-y), & \text{if } y>0,\\
	1, & \text{if } y\leq 0.
\end{cases}
\]
Therefore, splitting the integral at \(m=x\), we obtain
\begin{align*}
	\mathbb{P}\left(\bar{\theta}_{\text{full}}(\mathbf{U},2/3)>x\right)
	&=
	\int_{x/2}^{x}\exp\bigl(-(x-m)\bigr)f_M(m)\,\mathrm{d}m
	+
	\int_x^{\infty}f_M(m)\,\mathrm{d}m\\
	&=
	-\exp(-2x)+2\exp\left(-\frac{3}{2}x\right).
\end{align*}
Hence, the CDF is
\[
F_{\bar{\theta}_{\text{full}}(\mathbf{U},2/3)}(x)
=
1+\exp(-2x)-2\exp\left(-\frac{3}{2}x\right).
\]
Differentiating the CDF yields the PDF
\[
f_{\bar{\theta}_{\text{full}}(\mathbf{U},2/3)}(x)
=
-2\exp(-2x)+3\exp\left(-\frac{3}{2}x\right).
\]



\newpage


\section*{Code}
\addcontentsline{toc}{section}{Code}


